\chapter{Cod Sursă --- Runner LSTM}
\label{app:cod_lstm}

Această anexă include runner-ul LSTM separat, responsabil pentru consumul
cozii \texttt{epimodel:jobs:lstm}, antrenarea modelului recurent pe seria
\texttt{infected} și persistarea seriilor \texttt{observed}, \texttt{fitted}
și \texttt{predicted}. Spre deosebire de runner-ele compartimentale, această
implementare nu rezolvă ecuații diferențiale, ci antrenează o rețea LSTM
hand-rolled în NumPy (forward pass cu poartă de intrare, de uitare, de
ieșire și celulă; backward pass prin BPTT, cu Adam ca optimizator)
\textemdash{} detalii matematice în \cref{ch:lstm}. Split-ul cronologic
$80\%/20\%$ separă fereastra de antrenare de cea de test, normalizarea este
ajustată doar pe train pentru a evita scurgerea informației, iar metricile
MAE, RMSE și nRMSE sunt calculate pe fereastra de test. Validarea împotriva
predictorilor naivi (\cref{tab:lstm_baselines}) este executată separat prin
\texttt{scripts/lstm\_baselines.py}, nu de runner. Dependențele sunt
declarate în \texttt{code/runners/lstm/requirements.txt} (\texttt{numpy},
\texttt{redis}, \texttt{psycopg}), iar imaginea separată permite scalarea
independentă a calculului LSTM (costisitor, $\sim$$290$~s per rulare) față
de solverele ODE (sub-secunde).

\lstinputlisting[language=Python, caption={Runner-ul LSTM dedicat predicțiilor epidemiologice.}]{../code/runners/lstm/runner.py}
